This book teaches everything behind the \texttt{arbiter} project: what Jev is, how
calibration works, how the SDK is designed, and what the evidence says. It is a living
document, rebuilt as the project progresses (see the Build log in Part V).

\section*{Reading paths by goal}

\begin{description}[leftmargin=1.6cm,style=nextline]
  \item[\textcolor{clAccent}{10-minute tour}] Chapter~\ref{ch:building} (what we build) then
    Chapter~\ref{ch:jev} \S\ref{sec:jev-oneline} (what Jev does in one line) then
    Figure~\ref{fig:arch} (the architecture).
  \item[\textcolor{clAccent}{SDK user}] Part~II foundations, then Chapter~\ref{ch:arch}
    (architecture) and Chapter~\ref{ch:decisionpoints} (decision points).
  \item[\textcolor{clAccent}{Researcher}] Chapter~\ref{ch:calibration} (calibration),
    Chapter~\ref{ch:eval} (evaluating agents), and Part~IV (evidence).
  \item[\textcolor{clAccent}{Contributor}] Chapter~\ref{ch:engineering} (engineering
    practices) and the Build log (Part~V).
\end{description}

\section*{Conventions}

Every factual claim carries an evidence label so you always know how much to trust it:
\DOC{} verified in TypeSafe's documentation; \IND{} an independent measurement;
\VENDOR{} a vendor claim; \HYP{} our own hypothesis; \MEAS{} measured by us (with a run
ID). No \MEAS{} labels exist yet, because we have not run experiments.

Every figure and table caption ends with a \textbf{data label} in brackets stating where
its numbers come from: \emph{Illustrative} (a drawing, not data), \emph{Exact} (computed
from a formula), \emph{Reported by \dots} (someone else's measurement, cited), or
\emph{Measured by us}. This rule exists so that a teaching diagram is never mistaken for a
result.

\begin{keyidea}
The whole project rests on one honest distinction: a model's \emph{answer} and its
\emph{calibrated confidence in that answer} are two different things. Most of this book is
about the second one.
\end{keyidea}
